\chapter{АНАЛИТИЧЕСКИЙ ОБЗОР}

%%%%%%%%%%%%%%%%%%%%%%%%%%%%%%%%%%%%%%
%%% РОЛЬ АНИОНОВ В ПРИРОДНЫХ ВОДАХ %%%
%%%%%%%%%%%%%%%%%%%%%%%%%%%%%%%%%%%%%%

\section{РОЛЬ АНИОНОВ В ПРИРОДНЫХ ВОДАХ} % ~2стр.

Неорганические анионы растворённые в поверхностных и подземных водах являются основными компонентами, которые определяют качества этих вод, а именно: химический состав, буферную ёмкость и пригодность водоёма для питьевого, бытового, а также промышленного использования. В водной среде хлорид-ионы ($Cl^-$) представляют из себя достаточно подвижные и инертные анионы: они практически не адсорбируются на взвешенных частицах, не выпадают в осадок, кроме того хлор анион потребляется водной фауной в незначительно малых объёмах. Поэтому ион хлора(I) может выступать как важнейший индикатор галогенного загрязнения и протекания природных геологических процессов, связанных со смешением пресных континентальных и морских вод. Последнее поведение особенно характерно для прибрежных районов и районов, примыкающих к эстуариям рек. Так присутствие необычайно высоких концентраций хлоридов зачастую является сигналом проникновения солёной воды или загрязнение водоносных слоёв неочищенными городскими сточными водами. В незагрязнённых же водах хлорид-ионы находятся в резульате выщелачивания из магматических и осадочных пород, выпадения атмосферных осадков, а его содержания регулируется в соотвествие с активностью этих процессов \cite{pobozyApplicationCapillaryElectrophoresis2021b}. Экологическим воздействием избыточного содержания хлоридов является снижение биоразнообразия за счёт подавления активности восприимчивых видов насекомых, заместо которых распространяются инвазивные виды устойчивые к высокому содержанию солей. Для человека повышается риск развития сосудно-сердечных заболеваний и повышенного кровяного давления \cite{ kaushalFreshwaterSalinizationSyndrome2021}.

Сульфат-анионы ($SO_4^{2-}$) в природной воде выступают в качестве индикаторов интенсивности земплепользования, горнодобывающей деятельности и активности кислотных воздействий, в том числе совместно с соотношением хлорид анионов, сульфаты могут выступать как параметр для оценки риска коррозии свинцовых труб. Естественно сульфаты попадают в природные воды в результате выветривания осадочных пород, вулканической активности, растворения гипса и атмосферных осадков. Их повышенные содержания, вплоть до 250х кратного превышения по сравнению с незагрязнёнными водами, обеспечивается антропогенным воздействием: вымывание удобрение, слив неочищенных сточных вод, шахтные выбросы, а также сжиганием ископаемого топлива. Экологическим последствием повышенного содержания сульфатов является эвтрофикация в связи с высвобождением фосфора из осадочных пород, токсичность $H_2S$ для водных организов. Также отмечается воздействие на человеческий организм, а именно слабительное действие и раздраждение ЖКТ при употреблении такой воды. \cite{kaushalFreshwaterSalinizationSyndrome2021}

Ключевыми индикаторами состояния водных экосистем в том числе являются нитраты ($NO_3^-$) и нитриты ($NO_2^-$). Повышенные концентрации нитратов зачастую являются показателями антропогенного воздействия, в частности, интенсивном использовании сельскохозяственных удобрений и последующем их вымывании в поверхностные и грунтовые воды с помощью осадков. Кроме того, в отдельных случаях перелив из комбинированных канализационных систем в эпизоды сильных дождей может оказать сильное влияние на содержание концентраций азота. Мониторинг нитритов также может указать на микробные процессы, которые происходят в воде в анаэробных условиях в присутствии источников углеродов. Естественным источником соединений азота является круговорот этих соединений \cite{oliverImpactsClimateChange2024, arndtOnlineHighresolutionRealtime2026}. С точки зрения экологических последствий избыток нитратов и нитров в поверхностных водах ведёт к эвтрофикации, что в дальнейшем приводит к истощению кислорода с последующим высвобождением в воду цианотоксинов, а также сокращению биоразнообразия. В прибрежных зонах загрязнения ведёт к развитию у человека гипоксии или аноксии \cite{arndtOnlineHighresolutionRealtime2026}.

Общая щелочность является индикатором способности водоёма буферизировать поступающий $CO_2$, т.е. она отражает концентрацию гидрокарбонат ионов ($HCO_3^-$) и карбонат ионов ($CO_3^{2-}$) и некоторых других ионов, которые могут нейтрализовать свободные ионы водорода. Таким образом измерение pH и концентрации карбонат ионов служить индикатором загрязнения воды наравне с другими неорганическими анионами, за счёт того, что основным источником карбонатов в воде является поглощение антропогенного углекислого газа из-за сжигания углекислого топлива и сжигания лисов. Содержание карбонатов также определяется природными процессами, такими как осадки, выветривание минералов, эрозия и растворение геологических отложений, но в гораздо меньшей мере \cite{saalidongExaminingDynamicsRelationship2022}. После же повсеместного распространением промышленности по оценкам около 30\% всего выбрасываемого $CO_2$ поглощается только наземными водоёмами \cite{kellyNOAAOceanicAtmospheric2025}. В свою очередь чрезмерное содержание карбонат анионов в основном имеет последствия из-за повышения pH среды: негативным эффектом этого процесса является повышение жёсткости воды, снижение эффективности очистки воды. Особо высокие значения pH ($>9.3$) демонстрируют разрушающий эффект слизистых оболочек человека, глаз и кожи. \cite{saalidongExaminingDynamicsRelationship2022}. Напротив этому, для общего экологического состояния повышенные значения карбонатов являются благоприятным исходом за счёт противостояния глобальной проблеме закисления, снижение токсичности тяжёлых металлов и образования твёрдого $CaCO_3$ для внешних скелетов водных организмов \cite{saalidongExaminingDynamicsRelationship2022, figuerolaReviewMetaAnalysisPotential2021}.

%%%%%%%%%%%%%%%%%%%%%%%%%%%%%%%%%%%%%%%%%%%%%%
%%% СОВРЕМЕННЫЕ МЕТОДЫ ОПРЕДЕЛЕНИЯ АНИОНОВ %%%
%%%%%%%%%%%%%%%%%%%%%%%%%%%%%%%%%%%%%%%%%%%%%%

\section{СОВРЕМЕННЫЕ МЕТОДЫ ОПРЕДЕЛЕНИЯ АНИОНОВ} % ~3стр.

\textit{[ДОПОЛНИТЬ: обзор подходов (ионная хроматография, спектрофотометрия, потенциометрия, КЭ и др.), сравнение по чувствительности, селективности, стоимости, времени анализа].}

\textit{[ДОПОЛНИТЬ: по результатам обзора оформить таблицу сравнения методик и метрологических характеристик].}

%%%%%%%%%%%%%%%%%%%%%%%%%%%%%%%%%%%%%%%%
%%% МЕТОД КАПИЛЛЯРНОГО ЭЛЕКТРОФОРЕЗА %%%
%%%%%%%%%%%%%%%%%%%%%%%%%%%%%%%%%%%%%%%%

\section{МЕТОД КАПИЛЛЯРНОГО ЭЛЕКТРОФОРЕЗА} % ~5стр.

\subsection{История развития и основные принципы метода}

\textit{[ДОПОЛНИТЬ: краткая история КЭ и принцип разделения анионов].}

\subsection{Теоретические основы и устройство систем «Капель»}

\textit{[ДОПОЛНИТЬ: описание основных узлов прибора, принцип регистрации, ведущий электролит].}

\subsection{Применение КЭ для анализа объектов окружающей среды}

\textit{[ДОПОЛНИТЬ: примеры использования КЭ для природных вод, ссылки на источники].}

%%%%%%%%%%%%%%%%%%%%%%%
%%% ВЫВОДЫ ПО ГЛАВЕ %%%
%%%%%%%%%%%%%%%%%%%%%%%

\section{ВЫВОДЫ ПО ГЛАВЕ}

\textit{[ДОПОЛНИТЬ: краткое резюме и обоснование выбора методики].}
